\chapter{使用openai}
\label{openai}
\section{密钥保存}
\label{key_save}

\subsection{Windows设置方法}
\label{Windows设置环境方法}
\begin{figure}[htpb]  
    \centering  
    \begin{minipage}{0.8\textwidth}  
        \centering  
        \includegraphics[width=\textwidth]{part/自动化/openai/pic/高级系统设置.png}  
        \caption{环境变量}  
        \label{fig:环境变量}  
    \end{minipage}  
    \vspace{1em}  
    \begin{minipage}{0.8\textwidth}  
        \centering  
        \includegraphics[width=\textwidth]{part/自动化/openai/pic/新建系统变量.png}  
        \caption{新建系统变量}  
        \label{fig:新建系统变量}  
    \end{minipage}  
    \vspace{1em}  
    \begin{minipage}{0.8\textwidth}  
        \centering  
        \includegraphics[width=\textwidth]{part/自动化/openai/pic/设置apikey环境变量.png}  
        \caption{设置apikey环境变量}  
        \label{fig:设置apikey环境变量}  
    \end{minipage}  
\end{figure}



\subsection{Linux设置方法}
\label{Linux设置环境方法}
\begin{lstlisting}[style=demostyle,numbers=none]
vim ~/.bashrc
export OPENAI_API_KEY="sk-uzm6********************************949bE6"
source ~/.bashrc
\end{lstlisting}

\section{使用}
\label{openai使用}
最好的办法是将密钥保存在环境变量中，这样不会泄露密钥。
\begin{lstlisting}[style=demostyle,numbers=none]
from openai import OpenAI
client = OpenAI()# 读取环境变量中的密钥

response = client.chat.completions.create(
    model="gpt-3.5-turbo",
    messages=[
        {
            "role": "user",
            "content": "有限元是如何求解PDE的？"
        }
    ]
)
\end{lstlisting}
